\documentclass{article}

\usepackage{tabularx}
\usepackage{booktabs}
\usepackage{xstring}
\usepackage[round]{natbib}

\title{Reflection and Traceability Report on \progname}

\author{\authname}

\date{}

\input{../Comments}
%% Common Parts

\newcommand{\progname}{Project Proxi} % PUT YOUR PROGRAM NAME HERE
\newcommand{\authname}{Team \#23, Project Proxi
\\ Savinay Chhabra
\\ Amanbeer Singh Minhas
\\ Gourob Podder
\\ Ajay Singh Grewal} % AUTHOR NAMES                  

\usepackage{hyperref}
    \hypersetup{colorlinks=true, linkcolor=blue, citecolor=blue, filecolor=blue,
                urlcolor=blue, unicode=false}
    \urlstyle{same}
                                


\makeatletter
\let\OriginalUrl\url
\renewcommand{\url}[1]{%
\IfSubStr{#1}{/issues/}{%
\StrBehind{#1}{/issues/}[\IssueNumber]%
\href{#1}{\#\IssueNumber}%
}{%
\IfSubStr{#1}{/pull/}{%
\StrBehind{#1}{/pull/}[\PullNumber]%
\href{#1}{\#\PullNumber}%
}{%
\OriginalUrl{#1}%
}%
}%
}
\makeatother

\begin{document}

\maketitle

\section{Changes in Response to Feedback}

This section summarizes the changes made across project deliverables in
response to feedback from TAs, peer reviewers, and internal team review.
Each item identifies the feedback received, its source, and the specific
Revision 1 changes made in response. GitHub issue and pull request links
are included to support traceability to the corresponding tracked changes.

Peer review feedback was also considered during revision. When peer
feedback aligned with TA feedback or identified clear issues in clarity,
structure, traceability, or testability, the suggested changes were
incorporated. When peer feedback was ambiguous, conflicted with the
document's purpose, or introduced unnecessary implementation detail,
the team reviewed it but did not apply it. These decisions were made to
preserve consistency with project scope, black-box requirements
principles, and previously validated design choices.

\subsection{SRS and Hazard Analysis}

\subsubsection{SRS}

\noindent\rule{\textwidth}{0.6pt}

\noindent\textbf{Section 1.1.1 Start: SRS Feedback and Responses}

\noindent
The following items summarize TA feedback on the SRS and the specific
Revision 1 changes made in response.

\begin{itemize}

\item
\textbf{Feedback: Use symbolic constants instead of repeated numeric
values.} \\
\textbf{Source: TA Feedback} \\
\textbf{Changes Made:}
A dedicated symbolic constants subsection was added in Section 2.1.
Repeated thresholds for timing, accuracy, capacity, and related fit
criteria were moved there and replaced with named constants throughout
the document. This reduced duplication and made the same values usable
consistently across goals, requirements, and V\&V activities. \\
\textbf{Issue Link:}
\url{https://github.com/gkpodder/Proxi/issues/129} \\
\textbf{PR Link:}
\url{https://github.com/gkpodder/Proxi/pull/121}

\item
\textbf{Feedback: Sections should not start directly with subsections.} \\
\textbf{Source: TA Feedback} \\
\textbf{Changes Made:}
Introductory text was added to major sections so each section now
explains its purpose before listing subsections. This was applied across
the SRS so the document has clearer flow and better context for the
reader before detailed content begins. \\
\textbf{Issue Link:}
\url{https://github.com/gkpodder/Proxi/issues/129} \\
\textbf{PR Link:}
\url{https://github.com/gkpodder/Proxi/pull/121}

\item
\textbf{Feedback: Goals feel too low-level and resemble FRs or NFRs.} \\
\textbf{Source: TA Feedback} \\
\textbf{Changes Made:}
Section 1.2 was rewritten so the goals are high-level and stable rather
than implementation-like. Detailed measurable criteria were removed
from the goals and expressed later through functional and
non-functional requirements, with quantitative targets centralized in
the symbolic constants section. \\
\textbf{Issue Link:}
\url{https://github.com/gkpodder/Proxi/issues/129} \\
\textbf{PR Link:}
\url{https://github.com/gkpodder/Proxi/pull/121}

\item
\textbf{Feedback: Many spelling and wording issues.} \\
\textbf{Source: TA Feedback} \\
\textbf{Changes Made:}
The full SRS was reviewed for spelling, grammar, and wording problems.
This included correction of typos, general wording cleanup, and small
improvements to sentence flow so the document reads more consistently
and professionally. \\
\textbf{Issue Link:}
\url{https://github.com/gkpodder/Proxi/issues/129} \\
\textbf{PR Link:}
\url{https://github.com/gkpodder/Proxi/pull/121}

\item
\textbf{Feedback: Definitions and acronyms should be ordered
alphabetically.} \\
\textbf{Source: TA Feedback} \\
\textbf{Changes Made:}
The glossary and terminology material were reorganized alphabetically
to improve lookup and consistency. This makes terminology easier to
scan and aligns the document with expected reference-section style. \\
\textbf{Issue Link:}
\url{https://github.com/gkpodder/Proxi/issues/129} \\
\textbf{PR Link:}
\url{https://github.com/gkpodder/Proxi/pull/121}

\item
\textbf{Feedback: Tables need to be referenced in the document.} \\
\textbf{Source: TA Feedback} \\
\textbf{Changes Made:}
Tables were explicitly referenced in the surrounding text so their role
is explained rather than left implicit. This was done to improve
document coherence and to make the trace from discussion to tabular
evidence clearer for the reader. \\
\textbf{Issue Link:}
\url{https://github.com/gkpodder/Proxi/issues/129} \\
\textbf{PR Link:}
\url{https://github.com/gkpodder/Proxi/pull/121}

\item
\textbf{Feedback: Section 8.3 needed more spacing and readability.} \\
\textbf{Source: TA Feedback} \\
\textbf{Changes Made:}
Section 8.3 was reformatted for better spacing and readability so
individual product use cases are easier to follow. This change improved
visual separation between entries and reduced the sense of compressed
content in that section. \\
\textbf{Issue Link:}
\url{https://github.com/gkpodder/Proxi/issues/129} \\
\textbf{PR Link:}
\url{https://github.com/gkpodder/Proxi/pull/121}

\item
\textbf{Feedback: Most FRs were actually NFRs.} \\
\textbf{Source: TA Feedback} \\
\textbf{Changes Made:}
Requirements were revised to better distinguish functional behaviour
from non-functional quality attributes. Functional requirements were
kept focused on what the system does, while performance, usability,
safety, and related quality measures were expressed through the
appropriate non-functional sections and fit criteria. \\
\textbf{Issue Link:}
\url{https://github.com/gkpodder/Proxi/issues/129} \\
\textbf{PR Link:}
\url{https://github.com/gkpodder/Proxi/pull/121}

\item
\textbf{Feedback: Need a table or matrix for traceability.} \\
\textbf{Source: TA Feedback} \\
\textbf{Changes Made:}
Dedicated traceability from requirements to use cases was added so the
reader can follow how requirements relate to system behaviour. This
strengthened the bridge between the SRS and later V\&V artifacts by
making requirement coverage more explicit. \\
\textbf{Issue Link:}
\url{https://github.com/gkpodder/Proxi/issues/129} \\
\textbf{PR Link:}
\url{https://github.com/gkpodder/Proxi/pull/121}

\item
\textbf{Feedback: Many NFRs need fit criteria.} \\
\textbf{Source: TA Feedback} \\
\textbf{Changes Made:}
Measurable thresholds were clarified and made more consistent by using
the symbolic constants section and by tightening the wording of
non-functional criteria. This improved testability and reduced the
chance of inconsistent repeated numeric values across sections. \\
\textbf{Issue Link:}
\url{https://github.com/gkpodder/Proxi/issues/129} \\
\textbf{PR Link:}
\url{https://github.com/gkpodder/Proxi/pull/121}

\item
\textbf{Feedback: Constraints need rationale and clearer origin.} \\
\textbf{Source: TA Feedback} \\
\textbf{Changes Made:}
Rationale was added to key constraints so the document better explains
why they exist and whether they come from the problem context, user
needs, or team decisions. This improved transparency and reduced the
risk of presenting internal design choices as mandated constraints. \\
\textbf{Issue Link:}
\url{https://github.com/gkpodder/Proxi/issues/129} \\
\textbf{PR Link:}
\url{https://github.com/gkpodder/Proxi/pull/121}

\item
\textbf{Feedback: Personas are good, but their effect on requirements is
unclear.} \\
\textbf{Source: TA Feedback} \\
\textbf{Changes Made:}
The SRS now explicitly explains how personas affect requirement
priorities. This links user groups such as accessibility-focused users,
older adults, and power users to the emphasis placed on accessibility,
high-level interaction, safety, and ease of use requirements. \\
\textbf{Issue Link:}
\url{https://github.com/gkpodder/Proxi/issues/129} \\
\textbf{PR Link:}
\url{https://github.com/gkpodder/Proxi/pull/121}

\item
\textbf{Feedback: Assumptions should be labeled, including LLM-related
assumptions.} \\
\textbf{Source: TA Feedback} \\
\textbf{Changes Made:}
Assumptions were labeled explicitly as A1, A2, and so on in the
assumptions section. This includes assumptions about external services,
such as the use of off-the-shelf AI services, so those expectations are
visible and can later be examined through validation rather than being
left implicit. \\
\textbf{Issue Link:}
\url{https://github.com/gkpodder/Proxi/issues/129} \\
\textbf{PR Link:}
\url{https://github.com/gkpodder/Proxi/pull/121}

\item
\textbf{Feedback: External knowledge and references were not included.} \\
\textbf{Source: TA Feedback} \\
\textbf{Changes Made:}
External references were incorporated into the SRS so design and user
context claims are supported rather than left unreferenced. This helps
anchor the problem context and strengthens the professionalism of the
document. \\
\textbf{Issue Link:}
\url{https://github.com/gkpodder/Proxi/issues/129} \\
\textbf{PR Link:}
\url{https://github.com/gkpodder/Proxi/pull/121}

\item
\textbf{Feedback: No attempt to formalize any part of the document.} \\
\textbf{Source: TA Feedback} \\
\textbf{Changes Made:}
A formalized safety interaction protocol and a dedicated system state
diagram subsection were added. These additions introduce more precise
behavioural structure without turning the SRS into a design document. \\
\textbf{Issue Link:}
\url{https://github.com/gkpodder/Proxi/issues/129} \\
\textbf{PR Link:}
\url{https://github.com/gkpodder/Proxi/pull/121}

\item
\textbf{Feedback: Planning detail was thin.} \\
\textbf{Source: TA Feedback} \\
\textbf{Changes Made:}
Schedule and task-planning detail were expanded so the SRS more
clearly reflects phased development expectations and downstream
planning links. This helps connect requirements to the broader project
workflow. \\
\textbf{Issue Link:}
\url{https://github.com/gkpodder/Proxi/issues/129} \\
\textbf{PR Link:}
\url{https://github.com/gkpodder/Proxi/pull/121}

\end{itemize}

\subsubsection{Hazard Analysis}

\noindent\rule{\textwidth}{0.6pt}

\noindent\textbf{Section 1.1.2 Start: Hazard Analysis Feedback and
Responses}

\noindent
The following items summarize TA feedback on the Hazard Analysis and
the specific Revision 1 changes made in response.

\begin{itemize}

\item
\textbf{Feedback: Tables were not referenced in the document body.} \\
\textbf{Source: TA Feedback} \\
\textbf{Changes Made:}
Text was added to explicitly refer to the revision history table, the
main FMEA, and the roadmap tables. The introduction now points the
reader to Table 1, Section 5 introduces Table 2 directly, and Section 7
explains the role of Tables 3 and 4 in prioritizing mitigations. \\
\textbf{Issue Link:}
\url{https://github.com/gkpodder/Proxi/issues/138} \\
\textbf{PR Link:}
\url{https://github.com/gkpodder/Proxi/pull/125}

\item
\textbf{Feedback: Table formatting and readability issues.} \\
\textbf{Source: TA Feedback} \\
\textbf{Changes Made:}
The revision history was reformatted to one item per row, the ruled
tables were replaced with Booktabs styling, and the FMEA table was
rotated to improve readability and avoid compressed text. \\
\textbf{Issue Link:}
\url{https://github.com/gkpodder/Proxi/issues/138} \\
\textbf{PR Link:}
\url{https://github.com/gkpodder/Proxi/pull/125}

\item
\textbf{Feedback: Need privacy hazards for external LLM usage.} \\
\textbf{Source: TA Feedback} \\
\textbf{Changes Made:}
Privacy and external-service risk coverage was expanded. New attention
was given to what data may leave the system when external AI services
are used, and explicit assumptions A6 and A7 were added to state that
only minimum task-relevant data is sent and that users are informed
when external processing is required. Derived requirements such as
consent for external sharing and bounded logging were also included. \\
\textbf{Issue Link:}
\url{https://github.com/gkpodder/Proxi/issues/138} \\
\textbf{PR Link:}
\url{https://github.com/gkpodder/Proxi/pull/125}

\item
\textbf{Feedback: Hazard requirements and terminology should align with
the SRS and V\&V artifacts.} \\
\textbf{Source: TA Feedback} \\
\textbf{Changes Made:}
Terminology and identifiers were aligned so hazard-derived requirements
use names and references that connect more cleanly to the SRS and
validation documents. The FMEA also includes SRS references beside
hazards to strengthen cross-document traceability. \\
\textbf{Issue Link:}
\url{https://github.com/gkpodder/Proxi/issues/138} \\
\textbf{PR Link:}
\url{https://github.com/gkpodder/Proxi/pull/125}

\item
\textbf{Feedback: Need clearer system boundaries and assumptions.} \\
\textbf{Source: TA Feedback} \\
\textbf{Changes Made:}
System boundaries were clarified in Section 3 by separating internal
components from external systems such as the operating system,
third-party applications, and cloud services. Assumptions were stated
explicitly in Section 4, with extra emphasis on the new privacy-related
assumptions tied to external processing. \\
\textbf{Issue Link:}
\url{https://github.com/gkpodder/Proxi/issues/138} \\
\textbf{PR Link:}
\url{https://github.com/gkpodder/Proxi/pull/125}

\end{itemize}

\subsection{VnV Plan and Report}

\subsubsection{V\&V Plan}

\noindent\rule{\textwidth}{0.6pt}

\noindent\textbf{Section 1.3.1 Start: V\&V Plan Feedback and Responses}

\noindent
The following items summarize TA feedback on the V\&V Plan and the
specific Revision 1 changes made in response.

\begin{itemize}

\item
\textbf{Feedback: Use symbolic constants instead of repeated numeric
values.} \\
\textbf{Source: TA Feedback} \\
\textbf{Changes Made:}
Repeated numeric thresholds were replaced with symbolic constants and
the plan now includes a symbolic parameters appendix. This makes the
validation targets more consistent with the SRS and easier to update in
one place. \\
\textbf{Issue Link:}
\url{https://github.com/gkpodder/Proxi/issues/131} \\
\textbf{PR Link:}
\url{https://github.com/gkpodder/Proxi/pull/126}

\item
\textbf{Feedback: Many spelling mistakes, grammar issues, and weak flow.} \\
\textbf{Source: TA Feedback} \\
\textbf{Changes Made:}
The document was reviewed for spelling, grammar, and sentence flow.
Terminology and phrasing were cleaned up so the plan reads more
smoothly and professionally. \\
\textbf{Issue Link:}
\url{https://github.com/gkpodder/Proxi/issues/131} \\
\textbf{PR Link:}
\url{https://github.com/gkpodder/Proxi/pull/126}

\item
\textbf{Feedback: Missing text introducing Section 1.} \\
\textbf{Source: TA Feedback} \\
\textbf{Changes Made:}
An opening description was added before Section 1 so the document now
states what the plan covers, how evidence will be collected, and how it
traces back to the SRS. This gives the section a clearer entry point. \\
\textbf{Issue Link:}
\url{https://github.com/gkpodder/Proxi/issues/131} \\
\textbf{PR Link:}
\url{https://github.com/gkpodder/Proxi/pull/126}

\item
\textbf{Feedback: Proper stylization is macOS.} \\
\textbf{Source: TA Feedback} \\
\textbf{Changes Made:}
Platform terminology was standardized, including the proper stylization
of macOS, so the document uses consistent technical naming. \\
\textbf{Issue Link:}
\url{https://github.com/gkpodder/Proxi/issues/131} \\
\textbf{PR Link:}
\url{https://github.com/gkpodder/Proxi/pull/126}

\item
\textbf{Feedback: Table titles and table references were incorrect or
missing.} \\
\textbf{Source: TA Feedback} \\
\textbf{Changes Made:}
Table captions were positioned above tables and the text now refers to
tables directly, such as Table 1 for team roles. This improves
navigation and makes tabular evidence part of the argument rather than
isolated content. \\
\textbf{Issue Link:}
\url{https://github.com/gkpodder/Proxi/issues/131} \\
\textbf{PR Link:}
\url{https://github.com/gkpodder/Proxi/pull/126}

\item
\textbf{Feedback: No plan for how roles may change if needed.} \\
\textbf{Source: TA Feedback} \\
\textbf{Changes Made:}
A dedicated role reassignment and contingency plan was added in
Section 2.2. It specifies how blocked work is identified, how backup
reviewers are maintained, how reassignment is tracked in GitHub, and
how traceability records are updated afterward. \\
\textbf{Issue Link:}
\url{https://github.com/gkpodder/Proxi/issues/131} \\
\textbf{PR Link:}
\url{https://github.com/gkpodder/Proxi/pull/126}

\item
\textbf{Feedback: Checklists were not atomic enough.} \\
\textbf{Source: TA Feedback} \\
\textbf{Changes Made:}
Broad checklist statements were broken into smaller, independently
checkable items. For example, review checklists now separate ideas
such as identifier uniqueness, direct language, measurable outcome,
and requirement consistency, rather than hiding several checks inside
one vague bullet. \\
\textbf{Issue Link:}
\url{https://github.com/gkpodder/Proxi/issues/131} \\
\textbf{PR Link:}
\url{https://github.com/gkpodder/Proxi/pull/126}

\item
\textbf{Feedback: ``Run some WCAG checks'' was too vague.} \\
\textbf{Source: TA Feedback} \\
\textbf{Changes Made:}
The accessibility verification plan was made concrete. The plan now
refers to specific checks such as keyboard navigation, screen reader
compatibility, and WCAG 2.2-related accessibility validation instead
of relying on a generic statement. \\
\textbf{Issue Link:}
\url{https://github.com/gkpodder/Proxi/issues/131} \\
\textbf{PR Link:}
\url{https://github.com/gkpodder/Proxi/pull/126}

\item
\textbf{Feedback: It was not clear whether traceability was complete.} \\
\textbf{Source: TA Feedback} \\
\textbf{Changes Made:}
The plan now states requirement-to-test traceability more explicitly
and includes a dedicated traceability section for functional and
non-functional system tests. This clarifies intended completeness of
coverage. \\
\textbf{Issue Link:}
\url{https://github.com/gkpodder/Proxi/issues/131} \\
\textbf{PR Link:}
\url{https://github.com/gkpodder/Proxi/pull/126}

\item
\textbf{Feedback: Walkthroughs, inspections, and reviews lacked
specific plans.} \\
\textbf{Source: TA Feedback} \\
\textbf{Changes Made:}
Review activities were expanded into clearer planned stages, including
internal review, peer team review, walkthrough meetings, and
traceability checks. The plan now explains what each activity is for
and how issues will be recorded. \\
\textbf{Issue Link:}
\url{https://github.com/gkpodder/Proxi/issues/131} \\
\textbf{PR Link:}
\url{https://github.com/gkpodder/Proxi/pull/126}

\item
\textbf{Feedback: Chosen extras were not justified well enough.} \\
\textbf{Source: TA Feedback} \\
\textbf{Changes Made:}
The challenge/extras section now explains why the user manual and
usability report matter for this project specifically. Their value is
tied to onboarding, accessibility support, and iterative improvement
based on observed usability findings. \\
\textbf{Issue Link:}
\url{https://github.com/gkpodder/Proxi/issues/131} \\
\textbf{PR Link:}
\url{https://github.com/gkpodder/Proxi/pull/126}

\item
\textbf{Feedback: Stakeholder review session and survey goals were not
clear enough.} \\
\textbf{Source: TA Feedback} \\
\textbf{Changes Made:}
The plan was expanded to explain how stakeholder feedback and survey
results will be used, and the appendix now includes a usability survey
with acceptance criteria. This makes the intended outputs of those
activities more explicit and measurable. \\
\textbf{Issue Link:}
\url{https://github.com/gkpodder/Proxi/issues/131} \\
\textbf{PR Link:}
\url{https://github.com/gkpodder/Proxi/pull/126}

\end{itemize}

\subsubsection{V\&V Report}

\noindent\rule{\textwidth}{0.6pt}

\noindent\textbf{Section 1.3.2 Start: V\&V Report Feedback and Responses}

\noindent
The following items summarize TA feedback on the V\&V Report and the
specific Revision 1 changes made in response.

\begin{itemize}

\item
\textbf{Feedback: Some FR tests were really more like NFR tests.} \\
\textbf{Source: TA Feedback} \\
\textbf{Changes Made:}
The report now separates behavioural verification from timing and
latency evidence. Functional sections focus on whether required actions
occur, while timing, latency, throughput, and related measurements are
reported in the non-functional performance section. \\
\textbf{Issue Link:}
\url{https://github.com/gkpodder/Proxi/issues/136} \\
\textbf{PR Link:}
\url{https://github.com/gkpodder/Proxi/pull/127}

\item
\textbf{Feedback: The project name should be updated.} \\
\textbf{Source: TA Feedback} \\
\textbf{Changes Made:}
The document title and content were updated to refer consistently to
Project Proxi instead of a generic software engineering label. This
brings the report in line with the other project artifacts. \\
\textbf{Issue Link:}
\url{https://github.com/gkpodder/Proxi/issues/136} \\
\textbf{PR Link:}
\url{https://github.com/gkpodder/Proxi/pull/127}

\item
\textbf{Feedback: Performance results should include data, not only
means.} \\
\textbf{Source: TA Feedback} \\
\textbf{Changes Made:}
Rawer performance evidence was added through dedicated tables for quiet
and moderate-noise trials, common action completion times, and CPU
samples during extended runs. This makes the performance discussion more
transparent and easier to inspect. \\
\textbf{Issue Link:}
\url{https://github.com/gkpodder/Proxi/issues/136} \\
\textbf{PR Link:}
\url{https://github.com/gkpodder/Proxi/pull/127}

\item
\textbf{Feedback: Usability has good data so far, but more complete
results are expected elsewhere.} \\
\textbf{Source: TA Feedback} \\
\textbf{Changes Made:}
The report keeps the current usability-related evidence that supports
requirement coverage, while making it clear that fuller usability
analysis belongs in the dedicated usability report. This keeps the
V\&V Report focused while preserving consistency with the chosen extra. \\
\textbf{Issue Link:}
\url{https://github.com/gkpodder/Proxi/issues/136} \\
\textbf{PR Link:}
\url{https://github.com/gkpodder/Proxi/pull/127}

\item
\textbf{Feedback: Participant-count discussion should reflect capstone
constraints.} \\
\textbf{Source: TA Feedback} \\
\textbf{Changes Made:}
Usability-study wording was adjusted so it reflects realistic
small-scale evaluation within capstone limits, while still presenting
meaningful participant-based evidence and observed issues. \\
\textbf{Issue Link:}
\url{https://github.com/gkpodder/Proxi/issues/136} \\
\textbf{PR Link:}
\url{https://github.com/gkpodder/Proxi/pull/127}

\item
\textbf{Feedback: Traceability in the final report should be clearer.} \\
\textbf{Source: TA Feedback} \\
\textbf{Changes Made:}
The report includes requirement traceability matrices for functional
and non-functional requirements, as well as module-level unit-test
traceability matrices. This makes the relation between planned and
completed verification more explicit. \\
\textbf{Issue Link:}
\url{https://github.com/gkpodder/Proxi/issues/136} \\
\textbf{PR Link:}
\url{https://github.com/gkpodder/Proxi/pull/127}

\end{itemize}

\subsection{Design Documentation}

\subsubsection{Module Guide (MG)}

\noindent\rule{\textwidth}{0.6pt}

\noindent\textbf{Section 1.3.1 Start: Module Guide (MG) Feedback and Responses}

\noindent
The following items summarize TA feedback on the MG and the specific
Revision 1 changes made in response.

\begin{itemize}

\item
\textbf{Feedback: I don't know if the Types module is really SW
decision, since it is related to your requirements and not just generic
software engineering things like searching/sorting algorithms.} \\
\textbf{Source: TA Feedback} \\
\textbf{Changes Made:}
The MG was revised so \texttt{SD-Types} is no longer framed as a generic
software-decision module. Instead, it is presented as a shared type
module that supports the actual domain model used by Proxi. \\
\textbf{Issue Link:}
\url{https://github.com/gkpodder/Proxi/issues/132} \\
\textbf{PR Link:}
\url{https://github.com/gkpodder/Proxi/pull/128}

\item
\textbf{Feedback: Core data types is not an abstract data type.} \\
\textbf{Source: TA Feedback} \\
\textbf{Changes Made:}
The description of the core structures was updated so they are treated
as shared type and record definitions rather than being described as an
abstract data type. This makes the MG classification more accurate. \\
\textbf{Issue Link:}
\url{https://github.com/gkpodder/Proxi/issues/132} \\
\textbf{PR Link:}
\url{https://github.com/gkpodder/Proxi/pull/128}

\item
\textbf{Feedback: One table reference didn't work.} \\
\textbf{Source: TA Feedback} \\
\textbf{Changes Made:}
The broken table reference was corrected so the architecture discussion
now points readers to the intended content without a failed reference. \\
\textbf{Issue Link:}
\url{https://github.com/gkpodder/Proxi/issues/132} \\
\textbf{PR Link:}
\url{https://github.com/gkpodder/Proxi/pull/128}

\item
\textbf{Feedback: Some tables extend and overlap with page numbers
(LaTeX's \texttt{longtable} package should help with this).} \\
\textbf{Source: TA Feedback} \\
\textbf{Changes Made:}
The large traceability tables were converted to \texttt{longtable} so
that they can span pages cleanly without colliding with footer content
or page numbers. \\
\textbf{Issue Link:}
\url{https://github.com/gkpodder/Proxi/issues/132} \\
\textbf{PR Link:}
\url{https://github.com/gkpodder/Proxi/pull/128}

\item
\textbf{Feedback: Need to change the variable for your team name, so it
doesn't just say ``Software Engineering''.} \\
\textbf{Source: TA Feedback} \\
\textbf{Changes Made:}
The shared project and team macros were checked and used consistently so
these design documents refer to Project Proxi and the actual team rather
than generic placeholder text. \\
\textbf{Issue Link:}
\url{https://github.com/gkpodder/Proxi/issues/132} \\
\textbf{PR Link:}
\url{https://github.com/gkpodder/Proxi/pull/128}

\item
\textbf{Feedback: I find it a bit weird that the hyperlinks point to
the list of modules and not to the sections with more details about
them.} \\
\textbf{Source: TA Feedback} \\
\textbf{Changes Made:}
The MG navigation and cross-reference wording were cleaned up so it is
clearer how readers should move from the high-level hierarchy to the
later detailed module descriptions. \\
\textbf{Issue Link:}
\url{https://github.com/gkpodder/Proxi/issues/132} \\
\textbf{PR Link:}
\url{https://github.com/gkpodder/Proxi/pull/128}

\item
\textbf{Feedback: Should quickly describe the tables in the body of the
text.} \\
\textbf{Source: TA Feedback} \\
\textbf{Changes Made:}
Short explanatory text was added around the traceability tables so the
reader is told what each table shows and how it supports the design
argument. \\
\textbf{Issue Link:}
\url{https://github.com/gkpodder/Proxi/issues/132} \\
\textbf{PR Link:}
\url{https://github.com/gkpodder/Proxi/pull/128}

\item
\textbf{Feedback: Uses hierarchy diagram looks great! I love the
different shapes for different types of modules --- I would add a legend
for the different shapes.} \\
\textbf{Source: TA Feedback} \\
\textbf{Changes Made:}
The uses-hierarchy explanation was expanded so the meaning of the
various node shapes is described directly in the surrounding text,
providing the equivalent of a legend for the diagram. \\
\textbf{Issue Link:}
\url{https://github.com/gkpodder/Proxi/issues/132} \\
\textbf{PR Link:}
\url{https://github.com/gkpodder/Proxi/pull/128}

\item
\textbf{Feedback: Figure 2 is not referenced in the body of the text.} \\
\textbf{Source: TA Feedback} \\
\textbf{Changes Made:}
The MG now explicitly references the main UI figure in the surrounding
body text so it is introduced as part of the discussion rather than left
unexplained. \\
\textbf{Issue Link:}
\url{https://github.com/gkpodder/Proxi/issues/132} \\
\textbf{PR Link:}
\url{https://github.com/gkpodder/Proxi/pull/128}

\item
\textbf{Feedback: Modules in MG and MIS do not match (e.g. M10 is
 different).} \\
\textbf{Source: TA Feedback} \\
\textbf{Changes Made:}
Module names and numbering were synchronized across the MG and MIS so
that the same module identifiers refer to the same components in both
design documents. \\
\textbf{Issue Link:}
\url{https://github.com/gkpodder/Proxi/issues/132} \\
\textbf{PR Link:}
\url{https://github.com/gkpodder/Proxi/pull/128}

\end{itemize}

\subsubsection{Module Interface Specification (MIS)}

\noindent\rule{\textwidth}{0.6pt}

\noindent\textbf{Section 1.3.2 Start: MIS Feedback and Responses}

\noindent
The following items summarize TA feedback on the MIS and the specific
Revision 1 changes made in response.

\begin{itemize}

\item
\textbf{Feedback: You don't need the ``if audio devices exist'' style
of guard in the semantics, since the exception already captures that
case.} \\
\textbf{Source: TA Feedback} \\
\textbf{Changes Made:}
Redundant success guards were removed from the access-program semantics
so the state transitions are simpler and rely on the defined exception
behaviour instead of duplicating that logic in prose. \\
\textbf{Issue Link:}
\url{https://github.com/gkpodder/Proxi/issues/132} \\
\textbf{PR Link:}
\url{https://github.com/gkpodder/Proxi/pull/128}

\item
\textbf{Feedback: Use \texttt{exc :=} for the exceptions.} \\
\textbf{Source: TA Feedback} \\
\textbf{Changes Made:}
Exception descriptions were rewritten using the expected
\texttt{exc :=} notation so the MIS matches the course formatting
conventions more closely. \\
\textbf{Issue Link:}
\url{https://github.com/gkpodder/Proxi/issues/132} \\
\textbf{PR Link:}
\url{https://github.com/gkpodder/Proxi/pull/128}

\item
\textbf{Feedback: Use \texttt{out :=} for outputs instead of
sentences.} \\
\textbf{Source: TA Feedback} \\
\textbf{Changes Made:}
Output descriptions were updated to use \texttt{out :=} consistently,
which makes the behaviour specifications shorter, clearer, and more
uniform across modules. \\
\textbf{Issue Link:}
\url{https://github.com/gkpodder/Proxi/issues/132} \\
\textbf{PR Link:}
\url{https://github.com/gkpodder/Proxi/pull/128}

\item
\textbf{Feedback: Don't define types in the state variables. Add an
``Exported Types'' section in the syntax instead.} \\
\textbf{Source: TA Feedback} \\
\textbf{Changes Made:}
Type definitions were moved out of the state-variable descriptions and
into dedicated \texttt{Exported Types} subsections, which makes the
module interfaces easier to read and build from. \\
\textbf{Issue Link:}
\url{https://github.com/gkpodder/Proxi/issues/132} \\
\textbf{PR Link:}
\url{https://github.com/gkpodder/Proxi/pull/128}

\item
\textbf{Feedback: Some things aren't clear enough to fully build this.
For instance, \texttt{IntentType} is referenced throughout but does not
seem to be defined anywhere.} \\
\textbf{Source: TA Feedback} \\
\textbf{Changes Made:}
The missing shared definitions, including \texttt{IntentType}, were added
and tied together with the related records and parameter structures so
that the interfaces are much closer to ``enough to build.'' \\
\textbf{Issue Link:}
\url{https://github.com/gkpodder/Proxi/issues/132} \\
\textbf{PR Link:}
\url{https://github.com/gkpodder/Proxi/pull/128}

\item
\textbf{Feedback: Some parts of the document are somewhat formal, but no
one piece is fully formalized.} \\
\textbf{Source: TA Feedback} \\
\textbf{Changes Made:}
The \texttt{BH-Input} module was strengthened into a more fully
formalized state-based specification, with clearer invariants,
transitions, and local functions. \\
\textbf{Issue Link:}
\url{https://github.com/gkpodder/Proxi/issues/132} \\
\textbf{PR Link:}
\url{https://github.com/gkpodder/Proxi/pull/128}

\end{itemize}

\subsection{Problem Statement and Development Plan}

\subsubsection{Problem Statement and Goals}

\noindent\rule{\textwidth}{0.6pt}

\noindent\textbf{Section 1.4.1 Start: Problem Statement and Goals Feedback and Responses}

\noindent
The following items summarize TA feedback on the Problem Statement and
Goals document and the specific Revision 1 changes made in response.

\begin{itemize}

\item
\textbf{Feedback: I'd put the most important goals for the project at
the top, e.g. goal 3 is really what the project is all about and should
probably go first.} \\
\textbf{Source: TA Feedback} \\
\textbf{Changes Made:}
The goals section was revised to better emphasize the project's central
outcomes and make the main priorities more apparent to the reader. \\
\textbf{Issue Link:}
\url{https://github.com/gkpodder/Proxi/issues/137} \\
\textbf{PR Link:}
\url{https://github.com/gkpodder/Proxi/pull/128}

\item
\textbf{Feedback: Goal 5: average user rating.} \\
\textbf{Source: TA Feedback} \\
\textbf{Changes Made:}
The wording of the user-satisfaction goal was clarified so the expected
average rating and its accessibility-focused context are stated more
clearly. \\
\textbf{Issue Link:}
\url{https://github.com/gkpodder/Proxi/issues/137} \\
\textbf{PR Link:}
\url{https://github.com/gkpodder/Proxi/pull/128}

\item
\textbf{Feedback: Sections should always have text in them and not just
a subsection header right below.} \\
\textbf{Source: TA Feedback} \\
\textbf{Changes Made:}
Introductory text was added so major sections now provide context before
moving into their subsections, improving the document's flow and
readability. \\
\textbf{Issue Link:}
\url{https://github.com/gkpodder/Proxi/issues/137} \\
\textbf{PR Link:}
\url{https://github.com/gkpodder/Proxi/pull/128}

\end{itemize}

\subsubsection{Development Plan}

\noindent\rule{\textwidth}{0.6pt}

\noindent\textbf{Section 1.4.2 Start: Development Plan Feedback and Responses}

\noindent
The following items summarize TA feedback on the Development Plan and
the specific Revision 1 changes made in response.

\begin{itemize}

\item
\textbf{Feedback: What if the other roles have to change? Do you have a
plan for that?} \\
\textbf{Source: TA Feedback} \\
\textbf{Changes Made:}
A clear role-reassignment contingency plan was added to the Team Member
Roles section so the document explains how role changes are handled if
team availability or project needs shift. \\
\textbf{Issue Link:}
\url{https://github.com/gkpodder/Proxi/issues/137} \\
\textbf{PR Link:}
\url{https://github.com/gkpodder/Proxi/pull/128}

\item
\textbf{Feedback: Be specific about branches, branch naming standards,
etc.} \\
\textbf{Source: TA Feedback} \\
\textbf{Changes Made:}
The workflow plan was made more concrete around issue-driven work,
branch use, pull-request review, and GitHub-based progress tracking. \\
\textbf{Issue Link:}
\url{https://github.com/gkpodder/Proxi/issues/137} \\
\textbf{PR Link:}
\url{https://github.com/gkpodder/Proxi/pull/128}

\item
\textbf{Feedback: These risks are not the kind of feasibility risks we
are looking for in this part.} \\
\textbf{Source: TA Feedback} \\
\textbf{Changes Made:}
The proof-of-concept discussion was refocused away from broader concerns
and toward a true feasibility question tied to the capstone scope. \\
\textbf{Issue Link:}
\url{https://github.com/gkpodder/Proxi/issues/137} \\
\textbf{PR Link:}
\url{https://github.com/gkpodder/Proxi/pull/128}

\item
\textbf{Feedback: You should pick one specific thing and prove it is
possible to do.} \\
\textbf{Source: TA Feedback} \\
\textbf{Changes Made:}
The plan now focuses on demonstrating one concrete capability:
interpreting a natural-language request and successfully carrying out a
meaningful computer action. \\
\textbf{Issue Link:}
\url{https://github.com/gkpodder/Proxi/issues/137} \\
\textbf{PR Link:}
\url{https://github.com/gkpodder/Proxi/pull/128}

\item
\textbf{Feedback: Be very clear what you will demonstrate in November.} \\
\textbf{Source: TA Feedback} \\
\textbf{Changes Made:}
The November proof-of-concept demonstration was rewritten more clearly
as an end-to-end workflow showing Proxi receive a request, choose tools,
and complete a visible task. \\
\textbf{Issue Link:}
\url{https://github.com/gkpodder/Proxi/issues/137} \\
\textbf{PR Link:}
\url{https://github.com/gkpodder/Proxi/pull/128}

\item
\textbf{Feedback: Be sure to format/stylize words, especially company
names like GitHub, properly.} \\
\textbf{Source: TA Feedback} \\
\textbf{Changes Made:}
Naming and styling were cleaned up so terms such as GitHub are presented
consistently and professionally throughout the document. \\
\textbf{Issue Link:}
\url{https://github.com/gkpodder/Proxi/issues/137} \\
\textbf{PR Link:}
\url{https://github.com/gkpodder/Proxi/pull/128}

\item
\textbf{Feedback: In LaTeX, use `` and '' for proper quotations.} \\
\textbf{Source: TA Feedback} \\
\textbf{Changes Made:}
Quotation formatting was corrected so the document follows standard
LaTeX typographic conventions more closely. \\
\textbf{Issue Link:}
\url{https://github.com/gkpodder/Proxi/issues/137} \\
\textbf{PR Link:}
\url{https://github.com/gkpodder/Proxi/pull/128}

\item
\textbf{Feedback: Should be written in first-person from each member.
This feels summarized.} \\
\textbf{Source: TA Feedback} \\
\textbf{Changes Made:}
The reflection writing was revised to read more naturally in first
person and to better reflect individual team perspective. \\
\textbf{Issue Link:}
\url{https://github.com/gkpodder/Proxi/issues/137} \\
\textbf{PR Link:}
\url{https://github.com/gkpodder/Proxi/pull/128}

\item
\textbf{Feedback: Should include a link to the repo, not just the
license and project board.} \\
\textbf{Source: TA Feedback} \\
\textbf{Changes Made:}
A direct link to the main repository was added so readers can reach the
project codebase from the Development Plan without relying only on the
license or project-board links. \\
\textbf{Issue Link:}
\url{https://github.com/gkpodder/Proxi/issues/137} \\
\textbf{PR Link:}
\url{https://github.com/gkpodder/Proxi/pull/128}

\item
\textbf{Feedback: Need to justify the use of the various technologies
and explore their pros and cons.} \\
\textbf{Source: TA Feedback} \\
\textbf{Changes Made:}
The `Expected Technology` section was expanded so the selected tools are
now justified in terms of their fit for Proxi, along with their main
benefits, limitations, and tradeoffs. This makes the engineering-tool
choices more thoughtful and defensible rather than leaving them as a
simple list. \\
\textbf{Issue Link:}
\url{https://github.com/gkpodder/Proxi/issues/137} \\
\textbf{PR Link:}
\url{https://github.com/gkpodder/Proxi/pull/128}

\end{itemize}

\section{Extras}

Two approved extras were completed for this project: a user manual and
a usability testing study. Both were selected to improve onboarding,
accessibility support, and evaluation of real user interaction quality.

\subsection*{User Manual}

A detailed user manual was developed to support onboarding and effective
use of the system. The manual is available
\href{../Extras/UserManual/UserManual.pdf}{here}.

It includes:

\begin{itemize}
\item Setup and installation instructions for running the system.
\item Description of system features and supported interactions.
\item Step-by-step examples of common user tasks and workflows.
\item Guidance for accessibility-focused usage scenarios.
\end{itemize}

The manual ensures that new users, including those with limited
technical experience, can effectively understand and interact with the
system. It also reinforces consistency between intended system
behaviour and actual user interaction.

\subsection*{Usability Testing}

A usability testing study was conducted to evaluate system usability and
interaction effectiveness. The full report is available
\href{../Extras/UsabilityTesting/UsabilityTesting.pdf}{here}.

The study includes:

\begin{itemize}
\item Scenario-based tasks representing realistic user workflows.
\item Observation of user interactions and identification of issues.
\item Collection of user feedback through structured surveys.
\item Measurement of usability metrics such as task completion and
user satisfaction.
\end{itemize}

The results were used to identify usability issues and guide system
improvements. These findings also support validation of usability-
related non-functional requirements.
\section{Design Iteration (LO11 (PrototypeIterate))}

\plt{Explain how you arrived at your final design and implementation.  How did
the design evolve from the first version to the final version?} 

\plt{Don't just say what you changed, say why you changed it.  The needs of the
client should be part of the explanation.  For example, if you made changes in
response to usability testing, explain what the testing found and what changes
it led to.}

\section{Design Decisions (LO12)}

\plt{Reflect and justify your design decisions.  How did limitations,
 assumptions, and constraints influence your decisions?  Discuss each of these
 separately.}

\section{Economic Considerations (LO23)}

There is a growing market for AI-assisted productivity and accessibility
tools. Adoption of AI-powered systems has increased significantly in
recent years as users seek more efficient and intuitive ways to interact
with software systems~\citep{ai_adoption}. In particular, tools that
reduce reliance on traditional interfaces and support natural language
interaction are becoming increasingly relevant.

Proxi is currently developed as an open-source system with no immediate
plans for commercialization. However, a potential commercial version
could target professionals, accessibility-focused users, and enterprise
environments that benefit from automation and simplified workflows.

If commercialized, estimated development and maintenance costs would
primarily include infrastructure for AI services, system integration,
and ongoing updates. A reasonable estimate for transitioning to a
production-ready system would be approximately \$20,000--\$40,000 CAD,
depending on hosting requirements, API usage, and long-term support.

A possible pricing model could involve a subscription-based approach,
with pricing in the range of \$10--\$25 per user per month. This aligns
with existing subscription-based productivity tools, which commonly fall
within a similar pricing range~\citep{pricing}.

Assuming an average price of \$15 per month, the system would require
approximately 150--300 active users over one year to recover initial
costs. Beyond this point, additional users would contribute to profit,
making the model scalable.

Since Proxi is currently open-source, the primary focus is on attracting
users rather than generating revenue. User adoption can be increased by:

\begin{itemize}
\item Providing clear documentation and onboarding through the user manual.
\item Demonstrating system capabilities through example workflows.
\item Leveraging usability testing results to improve user experience.
\item Sharing the project through GitHub and academic networks.
\end{itemize}

The potential user base is large. Globally, a significant portion of the
population benefits from assistive or accessibility-focused technologies,
and demand for accessible computing solutions continues to grow
~\citep{accessibility}. In addition, students and professionals seeking
more efficient workflows represent a broad secondary user group.

Overall, while Proxi is not currently positioned as a commercial product,
the system demonstrates strong potential for future adoption and
scalability. These estimates are preliminary and intended to demonstrate
feasibility rather than provide exact financial projections.

\section{Reflection on Project Management (LO24)}

\plt{This question focuses on processes and tools used for project management.}

\subsection{How Does Your Project Management Compare to Your Development Plan}

\plt{Did you follow your Development plan, with respect to the team meeting plan, 
team communication plan, team member roles and workflow plan.  Did you use the 
technology you planned on using?}

\subsection{What Went Well?}

\plt{What went well for your project management in terms of processes and 
technology?}

\subsection{What Went Wrong?}

\plt{What went wrong in terms of processes and technology?}

\subsection{What Would you Do Differently Next Time?}

\plt{What will you do differently for your next project?}

\section{Ethics and Equity (LO18)}

Our team considered ethics and equity throughout the design and
development of Proxi, with a focus on accessibility, inclusivity, and
responsible use of AI-assisted technologies.

A key principle guiding the project was universal design. Proxi enables
users to interact with computer systems using natural language rather
than requiring technical knowledge of commands or interfaces. This
reduces barriers for users who may have limited technical experience,
including older adults and users with accessibility needs.

The system was designed to support equitable access by prioritizing
simplicity and clarity in interactions. For example, users are not
required to remember specific commands, and feedback is presented in a
way that is easy to understand. These design choices help ensure that
the system is usable by a broader range of users, rather than only those
with advanced technical skills.

Empathy played an important role in understanding stakeholder needs.
Personas representing accessibility-focused users and non-technical
users were used to guide requirements and design decisions. This helped
the team consider how different users might experience frustration,
confusion, or barriers when interacting with traditional systems, and
influenced the focus on intuitive interaction and error tolerance.

The team also considered ethical implications of using external AI
services. There is a potential risk that user data may be transmitted to
external systems. To address this, the system is designed to minimize
data sharing and ensure that only task-relevant information is used.
Users are also expected to be informed when external processing may
occur, supporting transparency and informed consent.

Despite these efforts, there are still limitations. The system relies on
off-the-shelf AI models, which may not perform equally well across all
users or contexts. This could lead to inconsistent experiences,
particularly for users with different speech patterns or accents. While
this limitation was recognized, it was not fully addressed within the
scope of the project.

Overall, the team made a conscious effort to design Proxi in a way that
promotes equitable access and usability, while also acknowledging areas
where further improvements are needed to better support all users.

\section{Reflection on Capstone}

\plt{This question focuses on what you learned during the course of the capstone project.}

\subsection{Which Courses Were Relevant}

\plt{Which of the courses you have taken were relevant for the capstone project?}

\subsection{Knowledge/Skills Outside of Courses}

\plt{What skills/knowledge did you need to acquire for your capstone project
that was outside of the courses you took?}

\bibliographystyle{plainnat}
\bibliography{../../refs/References}

\end{document}